\documentclass[aspectratio=169,10pt]{beamer}

% =========================================================
% PAQUETES SEGUROS
% =========================================================
\usepackage[spanish]{babel}
\usepackage[utf8]{inputenc}
\usepackage[T1]{fontenc}
\usepackage{amsmath,amssymb}
\usepackage{xcolor}

% =========================================================
% TEMA SIMPLE Y ESTABLE
% =========================================================
\usetheme{Madrid}
\setbeamertemplate{navigation symbols}{}

% =========================================================
% COLORES UST
% =========================================================
\definecolor{USTBlue}{RGB}{0,70,127}
\definecolor{USTGray}{RGB}{80,80,80}

\setbeamercolor{frametitle}{bg=USTBlue,fg=white}
\setbeamercolor{title}{bg=USTBlue,fg=white}
\setbeamercolor{structure}{fg=USTBlue}

% =========================================================
% COMANDOS
% =========================================================
\newcommand{\guion}[1]{
	\begin{block}{Guion}
		#1
	\end{block}
}

\newcommand{\resultado}[1]{
	\begin{alertblock}{Resultado}
		#1
	\end{alertblock}
}

% =========================================================
% DATOS
% =========================================================
\title{Guion Video Evaluación Unidad 1}
\subtitle{Pensamiento Matemático -- Psicología}
\author{Luis Tulio González}
\date{}

% =========================================================
\begin{document}
	
	% ---------------------------------------------------------
	\begin{frame}
		\titlepage
	\end{frame}
	
	% ---------------------------------------------------------
	\begin{frame}{Estructura del video}
		\begin{enumerate}
			\item Presentación
			\item Problema 1
			\item Problema 2
			\item Problema 3
			\item Problema 4
			\item Cierre
		\end{enumerate}
	\end{frame}
	
	% =========================================================
	\section{Problema 1}
	% =========================================================
	
	\begin{frame}{Problema 1}
		\[
		(p \to q)\land(\sim r \to s)
		\]
		
		\guion{
			Si la memoria a corto plazo tiene capacidad ilimitada, entonces la memoria sensorial dura varios minutos, y si la memoria a largo plazo no es prolongada, entonces el individuo usa la información para tomar decisiones.
		}
	\end{frame}
	
	% ---------------------------------------------------------
	\begin{frame}{Problema 1(b)}
		\[
		[(p\land \sim q)\lor r]\to s
		\]
		
		Valores:
		\[
		p=F,\quad q=F,\quad r=V,\quad s=V
		\]
		
		\[
		[(F\land V)\lor V]\to V
		\]
		
		\[
		V \to V = V
		\]
		
		\resultado{La proposición es verdadera}
	\end{frame}
	
	% =========================================================
	\section{Problema 2}
	% =========================================================
	
	\begin{frame}{Problema 2(a)}
		\[
		\sim p \to (q\lor \sim r)
		\]
		
		\guion{
			Si una persona no sigue las estrategias terapéuticas, entonces presenta dificultades emocionales o no mejora su bienestar.
		}
	\end{frame}
	
	% ---------------------------------------------------------
	\begin{frame}{Problema 2(b)}
		\[
		\sim p \to (q\lor \sim r)
		\]
		
		\guion{
			La tabla de verdad muestra valores verdaderos y falsos.
		}
		
		\resultado{Es una contingencia}
	\end{frame}
	
	% =========================================================
	\section{Problema 3}
	% =========================================================
	
	\begin{frame}{Problema 3}
		\[
		A=\{P1,P3,P6,P7,P9,P15,P17,P18\}
		\]
		
		\[
		|A|=8
		\]
		
		\resultado{Participan en exactamente dos intervenciones}
	\end{frame}
	
	% ---------------------------------------------------------
	\begin{frame}{Problema 3(c)}
		\[
		\{P3,P4,P8,P9,P12,P13,P16,P17,P20\}
		\]
		
		\guion{
			Personas que están solo en H o en R y C simultáneamente.
		}
	\end{frame}
	
	% =========================================================
	\section{Problema 4}
	% =========================================================
	
	\begin{frame}{Problema 4(a)}
		\[
		(A-E)\cup(S\cap E)
		\]
		
		\[
		= \{2,4,8,9,10,14,16,20\}
		\]
		
		\resultado{Resultado correcto}
	\end{frame}
	
	% ---------------------------------------------------------
	\begin{frame}{Problema 4(b)}
		\[
		(A-S)^c\cap A^c - E
		\]
		
		\[
		= \{1,5,7,11,13,17,19\}
		\]
		
		\[
		|A|=7
		\]
		
		\resultado{Cardinalidad correcta}
	\end{frame}
	
	% =========================================================
	\section{Cierre}
	% =========================================================
	
	\begin{frame}{Cierre}
		\guion{
			En esta evaluación se trabajó lógica proposicional, tablas de verdad, diagramas de Venn y operaciones con conjuntos. Cada resultado fue justificado paso a paso.
		}
	\end{frame}
	
\end{document}